\scanpage{236}

\mspage{- 188 -}

IX

\[
 \underline R_U/\underline R_U^f \simeq M_U \ , \tag*{(13.5.12)}
\]

et l'homomorphisme $i$ de (13.3.4) peut s'écrire comme un homomorphisme

\[
 i:M_U\longrightarrow A^{\handwritten{\natural}}\ , \tag*{(13.5.13)}
\]

qu'on retrouvera dans un contexte un peu différent au par. suivant.

14. \uline{Liens avec la théorie transcendante} : \uline{cas rigide-analytique}. \uline{L'homomorphisme canonique} $\underline M_K\longrightarrow A_K^{\handwritten{\natural}\unsure{o|0}}$.

14.1. Nous supposons connu ici les fondements de la théorie des espaces rigide analytiques, pour lesquels on pourra consulter [12] [13] [34].

Soit $S$ un trait complet, $A_K$ un schéma abélien sur son corps des fractions. On a construit alors dans 8.1, par descente à partir de l'extension de Raynaud dans le cas semi-stable, une extension $G_K=A_K^{\handwritten{\natural}o}$ (8.1.1)

\[
 0\longrightarrow T_K\longrightarrow G_K\longrightarrow B_K\longrightarrow0 \tag*{(14.1.1)}
\]

d'un schéma abélien $B_K$ par un tore $T_K$. Considérant le schéma abélien dual $A'_K$, on a de même une extension $G'_K=A_K^{\prime\handwritten{\natural}o}$

\[
 0\longrightarrow T'_K\longrightarrow G'_K\longrightarrow B'_K\longrightarrow0\ , \tag*{(14.1.2)}
\]

où $B'_K$ s'identifie d'ailleurs au schéma abélien dual de $B_K$ (8.2). On désigne, comme dans (8.1.8), par $\underline M_K$ et $\underline M'_K$ les réseaux sur $K$ duals de $T'_K$ et de $T_K$ respectivement :

\[
 \underline M_K=D(T'_K)\ ,\quad\underline M'_K=D(T_K)\ . \tag*{(14.1.3)}
\]

\handwritten{\unsure{506}}

\mspage{506}

% archive sheet number: 235
